\section{Rev C --- Industrial Demonstration Family}
\label{sec:rev-c}

Rev~C contains only circuits that already worked in Rev~B. No architectural experimentation happens on Rev~C.

\begin{diagram}
                    ALFRED REV C

            +------------+------------+
            |            |            |
            v            v            v

      E2B-Control   E2B-Light      Gateway

        LAN8660      LAN8661     MCU+LAN8651
\end{diagram}

Made visually and electrically recognizable as one product family (Section~\ref{sec:platform-strategy}). Improvements over Rev~B:

\begin{itemize}[itemsep=0.1em]
\item Locking connectors (power and T1S).
\item Mechanical mounting suited to a demo fixture/enclosure.
\item PCB outline consistency across the three boards.
\item Enclosure compatibility.
\item Labeling.
\item Cable management.
\item ESD/transient protection, re-verified after connector changes.
\item BOM quality (replace any bench-expedient part with the intended representative part).
\item Assembly/test procedure, documented and repeatable.
\item Configuration simplicity.
\end{itemize}

\begin{decision}[title={What Rev C is not}]
Rev~C does not claim automotive production qualification of any kind. It is customer-demo-quality engineering hardware -- credible enough for a serious customer conversation, not a production part number. If Rev~C threatens the November~30 deadline, the program demonstrates on validated Rev~B hardware instead (Section~\ref{sec:feature-cuts}); architecture validation is never sacrificed to produce prettier PCBs.
\end{decision}
